\clearpage
\begin{landscape}
\begingroup
\fontsize{10}{13}\selectfont
\fancyhead[L]{\small NACHLASS. STÖRUNGEN DER PALLAS}
\renewcommand{\headrulewidth}{0pt}
\pdfbookmark[2]{Allgemeine Störungen, Erste Rechnung, [13]}{pallas-V-fragment13}
\begin{center}[13.]\end{center}
\par\noindent\hspace*{\parindent}[Es handelt sich nun darum, die mittlern Elemente zu ermitteln, welche, mit den im Vorigen berechneten Störungen verbunden, die Bewegung der Pallas darstellen. Sie werden erhalten, wenn man aus \SourceLeaf{533}den Seite 483---488 mit Hülfe der speciellen Störungen hergeleiteten Resultaten die Elemente der Pallas für die beobachteten sieben Oppositionen herleitet und an diese die nach den vorstehenden Artikeln zu berechnenden Störungen mit umgekehrten Zeichen anbringt.
\par\noindent\hspace*{\parindent}Nach den Seite 486 gegebenen Störungswerthen und unter Benutzung der dort gültigen Normalelemente ergeben sich für die sieben Oppositionen die folgenden osculirenden Elemente:
\begingroup
\small\setlength{\tabcolsep}{5pt}\renewcommand{\arraystretch}{1.18}
\begin{center}
\begin{tabular}{c|r|r|r|r|r|r}
Opposition&\(n\)&\(\varepsilon\)&\(\varphi\)&\(\tilde\omega\)&\(i\)&\(\AscendingNode\)\\
I&\(770''{,}7240\)&\(254^0\,56'\,42''{,}87\)&\(14^0\,12'\,42''{,}7\)&\(121^0\;7'\,54''{,}2\)&\(34^0\,37'\,44''{,}3\)&\(172^0\,28'\,21''{,}3\)\\
II&\(3919\)&\(57\;17{,}42\)&\(11\;28{,}5\)&\(6\;38{,}8\)&\(33{,}1\)&\(38{,}8\)\\
III&\(1365\)&\(56\;33{,}37\)&\(10\;53{,}9\)&\(3\;18{,}2\)&\(9{,}7\)&\(21{,}7\)\\
IV&\(4995\)&\(56\;39{,}18\)&\(10\;5{,}4\)&\(0\;54{,}1\)&\(22{,}5\)&\(8{,}4\)\\
V&\(1737\)&\(59\;5{,}71\)&\(8\;5{,}7\)&\(120\;59\;34{,}5\)&\(28{,}0\)&\(27\;37{,}3\)\\
VI&\(768{,}2675\)&\(56\;48{,}96\)&\(1\;31{,}9\)&\(45\;48{,}4\)&\(36\;36{,}9\)&\(27\;34{,}8\)\\
VII&\(8375\)&\(53\;38{,}45\)&\(13\;58\;52{,}2\)&\(48\;30{,}7\)&\(35\;12{,}2\)&\(25\;37{,}0\)
\end{tabular}
\end{center}
\endgroup
\par\noindent\hspace*{\parindent}Hier soll nur die Ableitung der halben grossen Axe und der Excentricität als Beispiel gegeben werden.
\par\noindent\hspace*{\parindent}Indem man die folgenden Werthe der mittlern Anomalien für die sieben Oppositionen zum Grunde legt:
\begingroup
\small\setlength{\tabcolsep}{5pt}\renewcommand{\arraystretch}{1.18}
\begin{center}
\begin{tabular}{@{}lrrrrrrr@{}}
\(M'\)&\(177^0\;3'\,19''\)&\(212^0\,33'\;2''\)&\(250^0\,28'\,46''\)&\(293^0\,45'\,48''\)&\(331^0\;5'\,39''\)&\(6^0\,13'\,25''\)&\(49^0\,11'\,33''\)\\
\(M\)&\(139\;19\;46\)&\(230\;47\;3\)&\(328\;28\;13\)&\(80\;2\;26\)&\(176\;13\;39\)&\(266\;44\;47\)&\(17\;27\;34\)
\end{tabular}
\end{center}
\endgroup
\par\noindent ergeben sich aus den numerischen Werthen des Artikels [10] die Störungen der Logarithmen der halben grossen Axe, wie folgt:
\[
-5648\quad-4381\quad-3424\quad-4729\quad-3516\quad+3776\quad+1683.
\]
\par\noindent Nach den obigen Elementen ist nun
\[
\log a\quad 0{,}442\,0718\quad0{,}442\,1966\quad0{,}442\,2926\quad0{,}442\,1562\quad0{,}442\,2786\quad0{,}442\,9961\quad0{,}442\,7813
\]
\par\noindent also der mittlere Werth der halben grossen Axe:
\[
0{,}442\,6366\quad0{,}442\,6347\quad0{,}442\,6350\quad0{,}442\,6291\quad0{,}442\,6302\quad0{,}442\,6185\quad0{,}442\,6130
\]
\par\noindent also das Mittel
\[
\log a=0{,}442\,6282
\]
\par\noindent und
\[
n=769''{,}2443.
\]
\par\noindent\hspace*{\parindent}Für den Logarithmen des halben Parameters ergeben sich die Störungen:
\[
-10935\quad-8902\quad-7593\quad-8433\quad-5895\quad+5550\quad+5160.
\]
\par\noindent Die obigen Elemente geben \(p=a\cos\varphi^2\):
\[
\log p\quad0{,}415\,0729\quad0{,}415\,2768\quad0{,}415\,4097\quad0{,}415\,3249\quad0{,}415\,5744\quad0{,}416\,7077\quad0{,}416\,6607
\]
\par\noindent woraus der mittlere Werth von \(\log p\) folgt:
\[
0{,}416\,1664\quad0{,}416\,1670\quad0{,}416\,1690\quad0{,}416\,1682\quad0{,}416\,1639\quad0{,}416\,1527\quad0{,}416\,1447
\]
\par\noindent also im Mittel
\[
\log p=0{,}416\,1617.
\]
\SourceLeaf{534}
\par\noindent Zur Bestimmung der Excentricität erhält man aus den gefundenen Werthen von \(a\) und \(p\):
\begingroup
\small\setlength{\tabcolsep}{5pt}\renewcommand{\arraystretch}{1.18}
\begin{center}
\begin{tabular}{@{}lrrrrrrr@{}}
\(\log\cos\varphi^2\)&\(9{,}973\,5298\)&\(9{,}973\,5323\)&\(9{,}973\,5340\)&\(9{,}973\,5391\)&\(9{,}973\,5337\)&\(9{,}973\,5342\)&\(9{,}973\,5317\)\\
\(\varphi=\)&\(14^0\,4'\,24''{,}5\)&\(22''{,}2\)&\(20''{,}6\)&\(15''{,}7\)&\(20''{,}8\)&\(20''{,}4\)&\(22''{,}7\)
\end{tabular}
\end{center}
\endgroup
\par\noindent im Mittel wird:
\[
\varphi\;(1803)=14^0\,4'\,21''{,}0.
\]
\par\noindent Die jährliche Abnahme von \(\varphi\) beträgt \(11''{,}13\); also erhält man, bezogen auf das Jahr 1810:
\[
\varphi=14^0\,3'\,3''{,}1.]
\]
\begin{center}\rule{.24\linewidth}{.4pt}\end{center}
\clearpage
\endgroup
\end{landscape}
